\documentclass{article}

\usepackage{corl_2026} % Use this for the initial submission.
% \usepackage[final]{corl_2026} % Uncomment for the camera-ready ``final'' version.
% \usepackage[preprint]{corl_2026} % Uncomment for pre-prints (e.g., arxiv); This is like ``final'', but will remove the CORL footnote.

% Encoding
\usepackage[utf8]{inputenc}
\usepackage[T1]{fontenc}

% Core math
\usepackage{amsmath, amssymb, amsfonts}
\usepackage{mathtools}
\usepackage{bm}

% Figures and tables
\usepackage{graphicx}
\usepackage{booktabs}
\usepackage{tabularx}
\usepackage{multirow}
\usepackage{makecell}
\usepackage{colortbl}
\usepackage{siunitx}
\usepackage{rotating}
\usepackage{wrapfig}

% Formatting
\usepackage{microtype}
\usepackage{xcolor}
\usepackage{nicefrac}
% \usepackage{refcheck}

% Links (load once, near the end)
% \usepackage[bookmarks=true]{hyperref}

% Algorithms and code
\usepackage[ruled]{algorithm2e}
\usepackage{listings}

% Subfigures
\usepackage[caption=false,font=footnotesize]{subfig}

% Misc
\usepackage{multicol}
\usepackage{dsfont}
\usepackage{tabularray}

% ── Math macros (inline, no external macros.tex required) ─────────────────────
% \newcommand{\reals}{\mathbb{R}}
% \newcommand{\naturals}{\mathbb{N}}
% \newcommand{\expectation}{\mathbb{E}}
% \newcommand{\wspace}{\mathcal{W}}
% \newcommand{\trans}{\mathcal{T}}
% \newcommand{\simulate}{\textsc{Simulate}}
% \newcommand{\config}{\mathcal{C}}
% \newcommand{\states}{\mathcal{S}}
% \newcommand{\rvisit}{R_{\mathrm{visit}}}
% \newcommand{\rterm}{R_{\mathrm{term}}}
% \newcommand{\rpath}{R_{\mathrm{path}}}
% ─────────────────────────────────────────────────────────────────────────────

% Formatting and General
\newcommand{\etal}{\textit{et al. }}
\newcommand{\solution}{\textbf{Solution:}}
\newcommand{\qed}{\hfill \ensuremath{\Box}}
\definecolor{darkgreen}{rgb}{0.0, 0.5, 0.0} 

% Math Sets and Symbols
\newcommand{\reals}{\ensuremath{\mathbb{R}}}
\newcommand{\naturals}{\ensuremath{\mathbb{N}}}
\newcommand{\integers}{\ensuremath{\mathbb{Z}}}
\newcommand{\expectation}{\ensuremath{\mathbb{E}}}
\newcommand{\D}{\mathrm{d}}
\newcommand{\reg}{\ensuremath{\emph{R}}}

% Spaces and Sets
\newcommand{\config}{\mathcal{Q}}
\newcommand{\wspace}{\mathcal{W}}
\newcommand{\obst}{\mathcal{O}}
\newcommand{\goals}{\mathcal{G}}
\newcommand{\states}{\mathcal{S}}
\newcommand{\trans}{\ensuremath{\mathcal{T}}}

% Rewards
\newcommand{\rterm}{R_{\mathrm{term}}}
\newcommand{\rspray}{R_{\mathrm{spr}}}
\newcommand{\rpath}{R_{\mathrm{path}}}
\newcommand{\rvisit}{R_{\mathrm{visit}}}

% Theorem Environments
\newtheorem{problem}{Problem}
\newtheorem{remark}{Remark}
\newtheorem{definition}{Definition}
\newtheorem{proposition}{Proposition}
\newtheorem{lemma}{Lemma}
\newtheorem{theorem}{Theorem}

% Paired Delimiters
\DeclarePairedDelimiter\ceil{\lceil}{\rceil}
\DeclarePairedDelimiter\floor{\lfloor}{\rfloor}

% Operators
\DeclareMathOperator*{\argmax}{arg\,max}
\DeclareMathOperator{\simulate}{SIM}
\DeclareMathOperator{\init}{init}
\DeclareMathOperator{\free}{free}
\DeclareMathOperator{\new}{new}
\DeclareMathOperator{\pos}{POS}
\DeclareMathOperator{\goalf}{GS}
\DeclareMathOperator{\visited}{\mathbb{V}}

% Fallback definitions for general-formulation macros (no-op if already defined in macros.tex)
\providecommand{\obst}{\mathcal{O}}
\providecommand{\free}{\mathrm{free}}
\providecommand{\obs}{\mathrm{obs}}
\providecommand{\col}{\mathrm{col}}

% Fallback definitions for general-formulation macros (no-op if already defined in macros.tex)
\providecommand{\obst}{\mathcal{O}}
\providecommand{\free}{\mathrm{free}}
\providecommand{\obs}{\mathrm{obs}}
\providecommand{\col}{\mathrm{col}}

\title{Learning to Navigate in Stochatic and Uncertain Environments with Different Robot Dynamics}

% The \author macro works with any number of authors. There are two
% commands used to separate the names and addresses of multiple
% authors: \And and \AND.
%
% Using \And between authors leaves it to LaTeX to determine where to
% break the lines. Using \AND forces a line break at that point. So,
% if LaTeX puts 3 of 4 authors names on the first line, and the last
% on the second line, try using \AND instead of \And before the third
% author name.

% NOTE: authors will be visible only in the camera-ready and preprint versions (i.e., when using the option 'final' or 'preprint').
% 	For the initial submission the authors will be anonymized.

\author{
  Jane E.~Doe\\
  Department of Electrical Engineering and Computer Sciences\\
  University of California Berkeley
  United States\\
  \texttt{janedoe@berkeley.edu} \\
  %% examples of more authors
  %% \And
  %% Coauthor \\
  %% Affiliation \\
  %% Address \\
  %% \texttt{email} \\
  %% \AND
  %% Coauthor \\
  %% Affiliation \\
  %% Address \\
  %% \texttt{email} \\
  %% \And
  %% Coauthor \\
  %% Affiliation \\
  %% Address \\
  %% \texttt{email} \\
  %% \And
  %% Coauthor \\
  %% Affiliation \\
  %% Address \\
  %% \texttt{email} \\
}

\begin{document}
\maketitle
\thispagestyle{empty}
\pagestyle{empty}

%===============================================================================

\begin{abstract}
Autonomous multi-robot systems operating under real-world conditions must be robust to wind disturbances, actuator noise, sensor imperfections, and other environmental uncertainties. Yet most existing deep RL navigation benchmarks assume simplified, deterministic simulators with fixed robot dynamics and workspace geometry, limiting applicability to heterogeneous real-world deployment. In this paper, we present a formally specified framework for constructing environments for multi-robot navigation in stochastic, uncertain settings with heterogeneous robot dynamics. We formulate the cooperative target-coverage task as a constrained stochastic optimization problem in which a team of robots must visit all goal locations while minimizing path costs and avoiding collisions under stochastic wind disturbances, actuation noise, and observation uncertainty. We prove the equivalence of this optimization problem to a Partially Observable Markov Decision Process~(POMDP), providing a formal link between the POMDP reward and the constrained navigation objective that prior navigation RL benchmarks do not specify. We instantiate the general framework with two physically distinct platforms---unmanned aerial vehicles~(UAVs) with Newtonian point-mass dynamics and wheeled ground robots governed by non-holonomic bicycle kinematics---and release corresponding environments supporting stochastic wind, actuation noise, observation noise, and domain randomization~(DR). Both environments are directly compatible with standard deep reinforcement learning algorithms and support systematic ablation studies over reward structure, observation information, uncertainty sources, and DR conditions. We evaluate six state-of-the-art deep RL algorithms across ten environment configurations per platform, and conduct systematic ablation studies over reward terms, observation spaces, uncertainty modes, and domain randomization. We further quantify the simulation-to-real observation gap for the UAV platform in a CoppeliaSim robotics scene, establishing a baseline for transfer research.
\end{abstract}


\section{Introduction}

Autonomous multi-robot systems must reliably navigate complex, unstructured environments where stochastic disturbances---wind gusts, imprecise actuators, and noisy sensors---are unavoidable. Applications ranging from infrastructure inspection and environmental monitoring to search-and-rescue and agricultural field coverage all demand that robots coordinate efficiently under these real-world uncertainties \citep{dulacArnold2021, fan2020ijrr}. Moreover, different mission profiles call for fundamentally different robot platforms and motion models: aerial vehicles are often studied with simplified Newtonian point-mass dynamics in largely obstacle-free workspaces, whereas wheeled ground robots are governed by non-holonomic bicycle kinematics and must reason about static obstacles in cluttered environments. A policy or training framework that assumes a single set of motion dynamics or a single workspace geometry cannot transfer across these platforms without redesign.

Reinforcement learning has emerged as a compelling approach to multi-robot navigation because it can, in principle, learn robust closed-loop policies directly from interaction with a stochastic simulator \citep{chen2017icra, everett2018iros, fan2020ijrr}. However, the design of the RL environment comes with unique challenges \citep{reda2020learning, dulac2020empirical, jakub2023as, henderson2018deep, sutton2018reinforcement}. Most existing RL environments for multi-robot navigation make two simplifying assumptions that limit their real-world relevance: deterministic dynamics (no wind, no actuation or observation noise) and fixed workspace geometry. Policies trained in such settings are brittle when confronted with the continuous stochastic perturbations of physical deployment---a well-documented sim-to-real failure mode \citep{dulacArnold2021, tobin2017domain}. Furthermore, without a formal link between the environment's reward signal and the underlying navigation objective, it is difficult to interpret whether learned policies optimize the intended task \citep{reda2020learning, dulac2020empirical}.

This paper fills those gaps by introducing a dynamics-agnostic stochastic navigation framework and benchmark. We formulate the cooperative navigation task using a general robot model and a unified free-workspace representation that accommodates both obstacle-free and obstacle-cluttered settings. We then instantiate this framework with two Gymnasium-compatible environments---one for UAV teams and one for wheeled-robot teams---that share the same task structure, algorithmic pipeline, and experimental ablation protocol, yet expose completely different motion dynamics. We formulate the shared cooperative task as a constrained stochastic optimization problem, derive the equivalent POMDP, and prove that the optimal POMDP policy solves the optimization problem. Our specific contributions are:

\begin{enumerate}
    \item \textbf{Dynamics-agnostic framework.} We present a general constrained optimization formulation for multi-robot target coverage with collision avoidance, built on a platform-independent robot model~\citep{le2021icra} and a unified free-workspace representation that naturally supports workspaces with or without static obstacles. The framework accommodates arbitrary robot dynamics through the generic simulator map.

    \item \textbf{Optimization-to-POMDP equivalence.} We formally prove that the POMDP's optimal policy solves the constrained navigation optimization problem, providing theoretical grounding absent in prior works \citep{chen2017icra, everett2018iros, fan2020ijrr}.

    \item \textbf{Dual-platform stochastic navigation environments.} We present two novel environments as concrete instantiations of the general framework, enabling direct comparison across robot types under identical experimental conditions.

    \item \textbf{Comprehensive benchmarking.} We evaluate six deep RL algorithms---A2C \citep{mnih2016a2c}, ARS \citep{mania2018ars}, TRPO \citep{schulman2015trust}, PPO \citep{schulman2017ppo}, TQC \citep{kuznetsov2020tqc}, and CrossQ \citep{bhatt2024crossq}---across reward, observation, uncertainty, and domain-randomization ablations on both platforms, providing a rigorous evaluation suite.

    \item \textbf{Sim-to-real gap analysis.} We deploy trained UAV policies in a 3-D CoppeliaSim~\citep{coppeliaSim2013} scene and quantify performance degradation under inference-time sensor imperfections, establishing a baseline for transfer research.
\end{enumerate}

\section{Related Work}
\label{sec:recent}

\subsection{Navigation and Planning}
Over the past two decades, numerous methods have been developed for multi-robot navigation and control \citep{nafis2021, sam2005, mickey2021}. Classical planning approaches such as A* \citep{goto2003}, RRT \citep{rodriguez2006rrt}, RRT* \citep{karaman2011rrts}, PSO \citep{nafis2021}, and RA-MCPP \citep{mickey2021} rely on explicit modeling of path-planning heuristics or global optimization strategies. PDDL \citep{fox2003pddl21} and PDDL+ \citep{fox2006pddlplus} have been widely used for multi-robot planning with temporal constraints, with planners such as POPF \citep{coles2010popf} and OPTIC \citep{benton2012optic} handling continuous numeric effects. ScottyActivity \citep{fernandez2018scottyactivity} interleaves heuristic search with convex optimization for coordinated teams but restricts robot behaviors to first-order linear dynamics and plans offline without stochastic disturbances at execution time. Conflict-Based Search (CBS) \citep{sharon2015cbs} and extensions such as multi-constraint CBS for geometrically larger agents \citep{li2019cbscontinuous} provide optimal multi-agent path finding on discretized graphs with deterministic edge costs. Sampling-based methods with robot dynamics \citep{le2017icaps, le2019ral, le2021icra, bui2022iros} construct deterministic motion trees over precisely modeled dynamics. All of these methods produce plans susceptible to compounding errors under the continuous stochastic perturbations present in real environments. By optimizing expected cumulative reward under the full distribution of dynamics, RL offers a substantially more robust alternative for stochastic and uncertain environments.

\subsection{Reinforcement Learning}
Reinforcement Learning (RL) has emerged as a promising alternative, especially for problems requiring real-time decision-making and adaptation in high-dimensional spaces \citep{kober2014, chen2017icra, everett2018iros, bae2019as, tang2025, dulacArnold2021, ibarz2021ijrr}. Early RL-based navigation works \citep{kober2014, bae2019as} often considered discrete or simplified state-action spaces. RL-based approaches to multi-robot navigation have since progressed from grid-world settings to high-dimensional continuous environments \citep{chen2017icra, everett2018iros, long2017iral, fan2020ijrr, han2022iral}. \citet{chen2017icra} demonstrated decentralized collision avoidance using deep RL for pairs of robots, while \citet{everett2018iros} extended this to $n$-robot teams using LSTMs over observations of an arbitrary number of neighbors. \citet{fan2020ijrr} showed that distributed RL policies can coordinate large teams---including up to 100 robots in simulation---in cluttered environments. However, these works focus on collision avoidance or fixed-goal navigation and typically assume simplified, deterministic simulators: they do not model wind disturbances, actuation noise, or heterogeneous robot dynamics within a single benchmark. Our work extends this line of research to explicitly stochastic environments with a dynamics-agnostic framework instantiated on two physically distinct robot types, together with a formal guarantee linking the POMDP reward to the underlying constrained navigation objective.

Recent advances in multi-agent reinforcement learning (MARL) provide scalable frameworks for decentralized coordination under uncertainty. Value decomposition methods such as VDN \citep{sunehag2018vdn} and QMIX \citep{rashid2018qmix, rashid2020qmix} enable tractable learning by factorizing joint value functions, while actor--critic approaches such as MADDPG \citep{lowe2017maddpg} and MAPPO \citep{yu2022mappo} achieve strong empirical performance in cooperative settings. Off-policy algorithms such as TQC \citep{kuznetsov2020tqc} and CrossQ \citep{bhatt2024crossq} have demonstrated high sample efficiency in continuous control, making them natural candidates for robotic navigation tasks. Robust RL methods \citep{peng2018sim2real, tobin2017domain} use domain randomization (DR) to bridge the sim-to-real gap by training over a distribution of physical parameters. Our work builds on these directions by providing a benchmark that supports systematic evaluation of on-policy, off-policy, and gradient-free algorithms across two platforms and multiple uncertainty conditions.

\subsection{Environment Design}
The importance of carefully designed environments is well recognized \citep{reda2020learning, cobbe2020icml, dulac2020empirical, sutton2018reinforcement, henderson2018deep, jakub2023as}. Benchmarking platforms such as Google Research Football \citep{kurach2020google} and CityFlow \citep{zhang2019cityflow} provide standardized training environments but do not specify task-level optimality guarantees linking rewards to navigation constraints. Physics-based simulators such as MuJoCo \citep{todorov2012mujoco}, PyBullet \citep{coumans2021pybullet}, and Isaac Gym \citep{makoviychuk2021isaac} enable realistic modeling of robot dynamics and have been widely used in robotics control and sim-to-real transfer studies \citep{tassa2018deepmind, akkaya2019solving, tobin2017domain}. Building on a general robot model for multi-robot motion planning \citep{le2021icra}, our framework is dynamics-agnostic, explicitly incorporates multiple sources of physical uncertainty, and provides a formal proof linking the POMDP reward to the underlying constrained optimization objective.

\section{Preliminaries}

Let $\reals, \reals_{\geq 0}, \naturals$ be the set of real, non-negative real, and natural numbers, respectively. Let $\{x_i\}_{i=1}^n$ denote an ordered sequence of elements $x_1, \ldots, x_n$. The $L^2$ norm $\|.\|_2$ is simplified as $\|.\|$ for convenience. Let $\mathbf{I}_n$ denote the $n \times n$ identity matrix. Let $\mathbf{0}_n$ and $\mathbf{1}_n$ denote $n$-dimensional all-zeros and all-ones vectors. For $x, y \in \reals$, we define the intervals $[x,y] = \{ i \in \reals \mid x \leq i \leq y \}, [x,y) = \{ i \in \reals \mid x \leq i < y \}$ and $(x,y] = \{ i \in \reals \mid x < i \leq y \}$. Let $\mathds{1}_C$ be a boolean variable where $\mathds{1}_C=1$ if condition $C$ is true and $0$ otherwise. For a set of samples $S$, let $\Delta(S)$ denote probability distributions over $S$. For a continuous random variable $Y$ with probability density function $f$, the expected value (or expectation) is defined as $\expectation[Y] = \int_{-\infty}^\infty y f(y) dy$. A sample from distribution $\mathcal{X}$ with mean $\mu \in \reals$ and variance $\sigma^2 \in \reals_{\geq 0}$ is written as $x \sim \mathcal{X}(\mu, \sigma^2)$. Let $\mathcal{N}(\mu, \sigma^2)$ denote the univariate normal (Gaussian) distribution. Let $\mathcal{N}(\bm{\mu}, \mathbf{\Sigma})$ denote the multivariate Gaussian distribution with mean vector $\bm{\mu} \in \reals^n$ and covariance matrix $\mathbf{\Sigma} \in \reals^{n \times n}$, where $\mathbf{\Sigma}$ is symmetric positive semi-definite \citep{bishop2006pattern}. In RL, a policy $\pi: S \to \Delta (A)$ is a function that maps a state $s \in S$ to a probability distribution over the set of actions $A$. The action-value or Q-value of a policy $\pi$ at state $s$ and action $a$ is defined as the expected cumulative reward from following the policy after taking the action $a$ starting from $s$ and written as, $Q^{\pi}(s,a) = \expectation_{\pi}[\sum_{t\geq0}\gamma^t r_t | s_0=s, a_0=a]$ where $r_t \in \reals$ is the reward at time $t$ and $\gamma \in [0,1]$ is the discount factor. We use the notation and definitions from \citet{choset2005} to formulate the multi-robot navigation problem in the next section.


\section{Problem Formulation}
\label{sec:prob}
Consider a group of $n \in \naturals$ mobile robots $\mathcal{M} = \{ \mathcal{M}_1, \ldots, \mathcal{M}_n \}$ operating in a bounded $d$-dimensional workspace $\wspace \subseteq \reals^d$ that may contain a set of static obstacles $\obst$. Let $\wspace_{O_j}$ denote the region occupied by obstacle $j$. The obstacle-free (free) workspace is $\wspace^{\free} = \wspace \setminus \cup_j \wspace_{O_j}$. The robots cooperate to visit a set of $m \in \naturals$ target locations (or goals) $\mathcal{G} = \{ g_1, \ldots, g_m \}$. Let $p_j^g \in \wspace$ denote the position of target $j$ and $r_j > 0$ denote its coverage radius. Following \citet{le2021icra}, each robot $\mathcal{M}_i = \{ \mathcal{P}_i, \config_i, \states_i, s_i^{\init}, \mathcal{U}_i, f_i \}$ is described by its shape $\mathcal{P}_i$, configuration space $\config_i$, state space $\states_i$, initial state $s_i^{\init} \in \states_i$, control space $\mathcal{U}_i$, and dynamics function $f_i$. At configuration $q_i \in \config_i$, robot $i$ occupies region $\mathcal{P}_i(q_i) \subset \wspace$. The obstacle-occupied and collision configuration spaces are:
\begin{align}
    \config_i^{\obs} &= \{q_i \in \config_i \mid \exists  j : \mathcal{P}_i(q_i) \cap \wspace_{O_j} \neq \emptyset\}, \label{eq:obs_config} \\
    \config_i^{\col} &= \{q_i \in \config_i \mid \exists  j \neq i : \mathcal{P}_i(q_i) \cap \mathcal{P}_j(q_j) \neq \emptyset\} \cup \config_i^{\obs}. \label{eq:col_config}
\end{align}
Since the robots are moving, let $\config_i^{col}(t)$ denote the collision configuration for robot $i$ at time $t$. The free configuration space at time $t$ is, $\config_i^{\free}(t) = \config_i \setminus \config_i^{\col}(t)$. The state space $\states_i = \config_i \times \mathcal{V}_i$ augments the configuration space with motion parameters $\mathcal{V}_i$ (e.g., velocity, heading, steering angle). The dynamics of each robot satisfy,
\begin{equation}
    \dot{s}_i = f_i(s_i, u_i), \quad s_i \in \states_i,\ u_i \in \mathcal{U}_i. \label{eq:dynamics}
\end{equation}
The dynamics follow a forward Euler discretization of a continuous-time process \citep{feedback2020}. A new state $s_i^{\mathrm{new}} \in \states_i$ reached after a small time step $dt$ is computed by numerically integrating $f_i$:
\begin{equation}
    s_i^{\mathrm{new}} \gets \simulate(s_i, u_i, f_i, dt). \label{eq:simulate}
\end{equation}
Starting at $s_i^{\init}$ and applying a sequence of controls gives a dynamically-feasible trajectory $\zeta_i : [0,T] \to \states_i$. The environment experiences a stochastic wind with per-step magnitude $v_w(t) \sim \mathcal{N}(\bar{v}_w, \epsilon_w^2)$ and direction $\theta_w(t) \sim \mathcal{N}(\bar{\theta}_w, \epsilon_\theta^2)$, yielding wind vector
\begin{equation}
    \mathbf{w}(t) = v_w(t) [\cos\theta_w(t),  \sin\theta_w(t)]^T, \label{eq:wind}
\end{equation}
where $\bar{v}_w$, $\bar{\theta}_w$ are nominal values and $\epsilon_w^2$, $\epsilon_\theta^2$ are noise variances. How $\mathbf{w}(t)$ changes the dynamics depends on the platform and is specified in Section~\ref{sec:exp:dynamics}. Each target $g_j \in \goals$ carries a binary visit flag $b_j(t) \in \{0, 1\}$, initialized to $0$ and transitioning irreversibly to $1$ once any robot comes within $r_j$ of $p_j^g$:
\begin{equation}
    b_j(t+dt) = \max \Big( b_j(t),  \max_{i=1}^n  \mathds{1}_{\|p_i(t) - p_j^g\| \leq r_j} \Big). \label{eq:coverage}
\end{equation}
Let $\mathbf{b}(t) = (b_1(t), \ldots, b_m(t)) \in \{0,1\}^m$ denote the full visit flags at time $t$. Let $\mathcal{J} : \states_i \times \mathcal{U}_i \to \reals_{\geq 0}$ be a per-step cost function:
\begin{equation}
    \mathcal{J}(s_i(t), u_i(t)) = \alpha_1 \|u_i(t)\|^2 + \alpha_2 \|v_i(t)\| + \alpha_3, \label{eq:cost}
\end{equation}
where $\alpha_1 \|u_i\|^2$ is the energy cost, $\alpha_2 \|v_i\|$ is the movement cost, $\alpha_3$ is the per-step time cost, and $\alpha_1, \alpha_2, \alpha_3 \in \reals_{\geq 0}$. The objective is to generate $n$ trajectories $\zeta_1,\ldots.,\zeta_n$ for the robots by applying $n$ control signals $u_1,\ldots,u_n$ that solves the stochastic optimization problem:
\begin{flalign}
 & \min_{u_1, \ldots, u_n} \mathbb{E} \Big[ \sum_{t=0}^{T} \sum_{i=1}^{n} \mathcal{J}(\zeta_i(t), u_i(t))  \Big|  \zeta_i(0) = s_i^{\init}, \mathbf{b}(0) = \mathbf{0}_m \Big] \label{eq:opt} \\
\text{s.t.,} \nonumber \\
\text{(i) }  & \sum_{g_j \in \goals} b_j(T) = m, \label{eq:const1} \\
\text{(ii) } & \forall t \in [0,T], \forall \mathcal{M}_i \in \mathcal{M}: \mathcal{P}_i(q_i(t)) \subset \wspace^{\free}, \label{eq:const2} \\
\text{(iii) }& \forall t \in [0,T], \forall \mathcal{M}_i \neq \mathcal{M}_j: \mathcal{P}_i(q_i(t)) \cap \mathcal{P}_j(q_j(t)) = \emptyset. \label{eq:const3}
\end{flalign}
Constraint~(i) ensures all target locations are covered; constraint (ii) enforces all robots to be in the free workspace; and constraint (iii) prevents inter-robot collisions.

\section{Environment Design}
\label{sec:mdp}

The optimization problem is formulated as a POMDP $M = (S, A, \Omega, O, \trans, R, \gamma)$. The \textbf{State Space} $S = \prod_{i=1}^n \states_i \times \{0,1\}^m$ captures the joint state $s_t = \big(\{s_i(t)\}_{i=1}^n,  \mathbf{b}(t)\big)$, where $\states_i$ is the platform-specific state space defined in Section~\ref{sec:exp:dynamics}. The \textbf{Action Space} $A = \mathcal{U}_1 \times \cdots \times \mathcal{U}_n$, where $\mathcal{U}_i$ is the platform-specific control space. The \textbf{Observation Space} $\Omega$ provide a noisy, partial view of the true state. The visit history $\mathbf{b}(t)$ is compressed into a single scalar via the binary-to-decimal encoding $\phi(\mathbf{b}) = \sum_{j=1}^m b_j 2^{j-1}$. Gaussian noise $\mathcal{N}(\mathbf{0}_c, \sigma_\mathrm{obs}^2 \mathbf{I}_c)$ is added to all observations where $\sigma_\mathrm{obs}^2 \in \reals_{\geq 0}$ is the joint observation variance and $c \in \naturals$ is the observation dimension. Platform-specific observation dimensions for each mode are given in Table~\ref{tab:obs_modes} in Section~\ref{sec:exp:dynamics}. The \textbf{Observation Function} $O : S \to \Delta(\Omega)$ maps each true state to a distribution over noisy observations as described above. The \textbf{Transition Function} $\trans : S \times A \to \Delta(S)$ governs the stochastic dynamics:
\begin{equation}
    \trans(s_t, a_t) = \Big(\simulate(s_1, u_1, f_1, dt), \ldots, \simulate(s_n, u_n, f_n, dt),  b_{t+dt}\Big), \label{eq:transition}
\end{equation}
where $f_i$ incorporates wind and actuation noise, and the visit flags update via Eq.~\eqref{eq:coverage}. The \textbf{Reward Function} $R : S \times A \to \reals$ balances target coverage, collision avoidance, and path efficiency. Let $r_s, r_m, r_L \in \reals_{\geq 0}$ be some small, medium, and very large rewards ($r_s < r_m \ll r_L$). The total reward decomposes as:
\begin{alignat}{2}
    R(s_t, a_t) &= \rvisit(t) + \rterm(t) + \rpath(t), \label{eq:reward}\\
    \rvisit(t)  &= r_m \cdot |\{g_j \in \goals : b_j(t + dt) = 1,  b_j(t) = 0\}|, \label{eq:r_visit} \\
    \rterm(t)   &=
    \begin{cases}
        -r_L, & \text{if } \exists \mathcal{M}_i: q_i \in \config_i^{\col}(t) \\
        +r_L, & \text{else if } \sum_{g_j \in \goals} b_j(t) = m \\
        0,    & \text{otherwise}
    \end{cases} \label{eq:r_term} \\
    \rpath(t)   &= - \sum_{i=1}^n \mathcal{J}(s_i(t), u_i(t))
                   - r_m \sum_{i=1}^n \mathds{1}_{p_i(t) \notin \wspace^{\free}}, \label{eq:r_path}
\end{alignat}
where $\rvisit$ rewards newly visited targets, $\rterm$ is the terminal reward for visiting all locations or exiting on collision, and $\rpath$ penalizes energy, speed, path length, and out-of-free-workspace steps. Finally, the \textbf{Discount Factor} $\gamma \in [0,1)$ incentivizes rapid target coverage. We define the following theorem to describe the relation between the optimization problem and the developed POMDP:

\begin{theorem}
\label{thm:optimization}
In the limit $\gamma \to 1$, the optimal policy $\pi^*$ for the POMDP $M$ starting at state $s = (\{s_i^{\mathrm{init}}\}_{i=1}^n,  \mathbf{0}_m)$ generates trajectories $\{\zeta_i\}_{i=1}^n$ that solve the optimization problem in Eq.~\eqref{eq:opt} with probability close to~1. Conversely, any solution of Eq.~\eqref{eq:opt} induces an optimal policy $\pi^*$ for $M$ with probability close to~1.
\end{theorem}

\textbf{Proof.} The proof below conditions on the true state $s_{t}$, which is standard practice when establishing reward-shaping equivalences; it corresponds to analyzing the induced MDP on beliefs in the limit of a sufficient statistic. For a finite horizon $[0,T]$, the definition of optimal policy $\pi^*$ and optimal action-value function $Q^{\pi^*}$ \citep{sutton2018reinforcement} are as follows:
\begin{align}
    \pi^*(s_t) &= \argmax_a Q^{\pi^*}(s_t, a), \label{eq:opt_pi} \\
    Q^{\pi^*}(s_t, a_t) &= \max_\pi \expectation_\pi \Big[\sum_{\tau=0}^T \gamma^\tau R(s_\tau, a_\tau)  \mid s_t, a_t\Big] \label{eq:opt_policy} \\
    = \max_\pi &\Big( \expectation_\pi \big[ \sum_\tau \gamma^\tau \rvisit(\tau) \mid . \big] + \expectation_\pi \big[ \sum_\tau \gamma^\tau \rterm(\tau) \mid . \big] + \expectation_\pi \big[ \sum_\tau \gamma^\tau \rpath(\tau) \mid . \big] \Big) \tag*{[From Eq. \eqref{eq:reward} and properties of $\mathbb{E}$]}
\end{align}

We prove both directions of the equivalence as follows.

\textbf{(Direction $\Rightarrow$) $\pi^*$ solves the optimization problem:} We proceed in two lemmas.

\begin{lemma}
\label{lem:const}
Any optimal policy $\pi^*$ must satisfy constraints~\textnormal{(i)}, \textnormal{(ii)}, and~\textnormal{(iii)} of Eq.~\eqref{eq:opt} with probability close to one.
\end{lemma}

\textit{Proof.} Because actuation and position noises are unbounded, no policy can eliminate the possibility of a constraint violation with probability exactly~1. We show that reward shaping makes violations overwhelmingly suboptimal.

\paragraph{Constraint (iii) --- no collisions.} Suppose $\pi^*$ induces a collision at some time $t^*$; then $\rterm(t^*) = -r_L$. Choosing $r_L \gg \sum_{\tau=0}^T \gamma^\tau (\rvisit(\tau) + \rpath(\tau))$ makes $Q^\pi \approx -r_L$, which is the lower bound. No such policy can be optimal.

\paragraph{Constraint (ii) --- obstacle and workspace containment.} Any policy placing robot $i$ outside $\wspace^{\free}$ (i.e., outside the workspace boundary or into an obstacle region) incurs a per-step penalty $-r_m$ in $\rpath$. Since $\wspace^{\free} = \wspace \setminus \bigcup_j \wspace_{O_j}$, this uniformly penalizes both boundary violations and obstacle penetrations. Keeping $\mathcal{P}_i(q_i) \subset \wspace^{\free}$ sets this term to zero, maximizing $\rpath$ since $r_m > 0$. When $\obst = \emptyset$, the condition reduces to $p_i \in \wspace$.

\paragraph{Constraint (i) --- all targets visited.} Visiting all $m$ targets yields $\rterm = +r_L$. Since $r_L$ dominates all other accumulated rewards, maximizing $\rterm$ ensures $\sum_{g_j \in \goals} b_j(T) = m$ at termination, satisfying constraint~(i). \hfill$\square$

\begin{lemma}
\label{lem:cost}
Among all policies satisfying constraints~\textnormal{(i)--(iii)}, maximizing $Q^\pi$ is equivalent to minimizing the cost objective of Eq.~\eqref{eq:opt}.
\end{lemma}

\textit{Proof.} By Lemma~\ref{lem:const} and constraint~(i), $\rterm(\tau) = +r_L$ at the terminal step $T$ and $0$ otherwise, so $\expectation_\pi[\sum_\tau \gamma^\tau \rterm(\tau) \mid .] = \expectation_\pi[\gamma^T r_L \mid .] \approx r_L$ for $\gamma \approx 1$.

For the visit reward, each target $j$ is visited exactly once by definition, contributing $r_m$ at the step it is first reached. Therefore:
\begin{equation}
    \expectation_\pi \Big[\sum_{\tau=0}^T \gamma^\tau \rvisit(\tau) \mid .\Big]
    = r_m \sum_{j=1}^m \expectation[\gamma^{t_j} \mid .]
    \approx r_m m, \label{eq:visit_const}
\end{equation}
where $t_j$ is the step when target $j$ is visited, and the approximation holds for $\gamma \approx 1$. This quantity is approximately constant (policy-independent for constraint-satisfying policies). By constraint~(ii), $\mathcal{P}_i(q_i(\tau)) \subset \wspace^{\free}$ for all $i, \tau$, so $\mathds{1}_{p_i \notin \wspace^{\free}} = 0$. Hence:
\begin{equation}
    \expectation_\pi \Big[\sum_\tau \gamma^\tau \rpath(\tau) \mid . \Big]
    \approx -\expectation_\pi \Big[\sum_\tau \gamma^\tau \sum_{i=1}^n \mathcal{J}(s_i(\tau), u_i(\tau)) \mid . \Big], \label{eq:path_const}
\end{equation}
From Eq.~\eqref{eq:opt_policy}:
\begin{flalign}
Q^{\pi^*}(s_t, a_t) &\approx \max_\pi \Big( r_L + r_m m - \expectation_\pi \Big[\sum_\tau \gamma^\tau \sum_i \mathcal{J}(s_i(\tau), u_i(\tau))\Big] \Big) \nonumber \\
&= r_L + r_m m - \min_\pi \expectation_\pi \Big[\sum_\tau \gamma^\tau \sum_i \mathcal{J}(s_i(\tau), u_i(\tau))\Big]. \nonumber
\end{flalign}
Here,  $r_L + r_m m$ is approximately constant (policy-independent). Consequently, maximizing $Q^{\pi^*}$ is precisely equivalent to minimizing the expected discounted cost in Eq.~\eqref{eq:opt}, proving the lemma. Therefore, by Lemmas~\ref{lem:const} and~\ref{lem:cost}, $\pi^*$ generates trajectories $\{\zeta_i\}_{i=1}^n$ that solve Eq.~\eqref{eq:opt}.

\textbf{(Direction $\Leftarrow$) The optimization solution induces $\pi^*$:} Let $\{\zeta_i^*\}_{i=1}^n$ be an optimal solution to Eqs.~\eqref{eq:opt}--\eqref{eq:const3}. Define a trajectory-tracking policy $\hat{\pi}$ that at each $t$ selects $\hat{\pi}(s_t) = \langle u_1^*(t), \ldots, u_n^*(t) \rangle$ from the optimal trajectory. By Lemma~\ref{lem:const}, $\zeta_i^*$ satisfies all three constraints with probability close to~1. Hence $\hat{\pi}$ achieves the maximum $\rterm$ and $\rvisit$ contributions and zero boundary penalty in $\rpath$. Now suppose for contradiction that there exists another policy $\pi'$ with $Q^{\pi'} > Q^{\hat{\pi}}$. By Lemma~\ref{lem:const} and the definition of $\pi^*$, $\pi'$ must also satistfy all three constraints. By Lemma~\ref{lem:cost}, $Q^{\pi'} > Q^{\hat{\pi}}$ implies:
\[
\sum_{t=0}^T \sum_{i=1}^n \mathcal{J} \big(\zeta_i^{\pi'}(t), u_i^{\pi'}(t)\big)
< \sum_{t=0}^T \sum_{i=1}^n \mathcal{J} \big(\zeta_i^*(t), u_i^*(t)\big),
\]
contradicting the optimality of $\{\zeta_i^*\}$. Therefore $\hat{\pi} = \pi^*$. Hence, Theorem~\ref{thm:optimization} holds in both directions. \hfill$\square$

\section{Experiments}
\label{sec:exp}

We present a comprehensive experimental evaluation to measure the following: performance of RL methods on both platforms, sensitivity to design choices, out-of-distribution (OOD) robustness under domain shift, and the simulation-to-real observation gap for the UAV platform.

% ─────────────────────────────────────────────────────────────────────────────
\subsection{Robot Dynamics and Observation Spaces}
\label{sec:exp:dynamics}
% ─────────────────────────────────────────────────────────────────────────────

Experimental evaluations are conducted for two groups of robots with different dynamics, instantiating the general framework of Section~\ref{sec:prob} as follows:

\paragraph{UAV dynamics (holonomic, Newtonian).}
Each UAV $\mathcal{M}_i$ has state $s_i = (p_i, v_i) \in \reals^4$ where $p_i \in \wspace$ is its 2-D position and $v_i \in [v_{\min}, v_{\max}]^2$ is its velocity vector. The control input $u_i \in [-1,1]^2$ specifies normalized 2-D thrust to change its velocity. Wind acts as an aerodynamic force that modifies velocity each step:
\begin{align}
    v_i(t+dt) &= \min \Big( \max \Big( v_i(t) + \tfrac{\kappa \mu_i(t) + \epsilon_a}{W_i} dt,  v_{\min} \Big),  v_{\max}\Big), \label{eq:uav_vel} \\
    p_i(t+dt) &= p_i(t) + \big(v_i(t) + \mathbf{w}(t)\big) dt, \label{eq:uav_pos}
\end{align}
where $\kappa > 0$ is the thrust power coefficient, $W_i > 0$ is the UAV mass, $\epsilon_a \sim \mathcal{N}(\mathbf{0}_2, \sigma_a^2 \mathbf{I}_2)$ is the actuation noise, and $\min, \max$ functions are applied element-wise. The UAV workspace $\wspace$ is a convex polygon without obstacles ($\obst = \emptyset$, $\wspace^{\free} = \wspace$).

\paragraph{Wheeled dynamics (non-holonomic, bicycle model).}
Each wheeled robot $\mathcal{M}_i$ has state $s_i = (x_i, y_i, \theta_i, v_i, \delta_i) \in \reals^5$, where $(x_i, y_i)$ is its 2-D position, $\theta_i$ is its heading angle, $v_i \in [v_{\min}, v_{\max}]$ is its longitudinal speed, and $\delta_i \in [\delta_{\min}, \delta_{\max}]$ is the front steering angle. The control input $u_i = (a_i, \dot{\delta}_i) \in [-1,1]^2$ specifies a normalized acceleration $a_i$ and steering rate $\dot{\delta}_i$. Unlike the UAV, wind acts on the wheeled robot as a lateral position drift rather than a velocity perturbation---reflecting the fact that ground vehicles are less susceptible to aerodynamic forces but can still be displaced by strong crosswinds:
\begin{align}
    v_i(t+dt) &= \min \big( \max \big( v_i(t) + a_s a_i dt,  v_{\min} \big),  v_{\max}\big), \label{eq:whl_vel} \\
    \delta_i(t+dt) &= \min \big( \max \big( \delta_i(t) + \dot{\delta}_i,  \delta_{\min} \big),  \delta_{\max}\big), \label{eq:whl_steer} \\
    \theta_i(t+dt) &= \theta_i(t) + \tfrac{v_i(t) \tan(\delta_i(t))}{\ell_i} dt, \label{eq:whl_head} \\
    \begin{bmatrix} x_i(t+dt) \\ y_i(t+dt) \end{bmatrix} &=
    \begin{bmatrix} x_i(t) \\ y_i(t) \end{bmatrix}
    + \Big( v_i(t) \begin{bmatrix}\cos\theta_i(t)\\\sin\theta_i(t)\end{bmatrix} + \mathbf{w}(t) \Big) dt, \label{eq:whl_pos}
\end{align}
where $a_s > 0$ is an acceleration scaling factor and $\ell_i > 0$ is the robot's wheelbase. The wheeled workspace is a rectangle $[0, W]\times[0, H]$ containing static polygon obstacles $\obst$. Table~\ref{tab:obs_modes} summarizes the observation space composition for each platform and mode.

\begin{table}[h]
\centering
\footnotesize
\caption{\small Observation space composition for each platform and mode ($N$~=~number of robots). UAV kinematics: position~$(2N)$ + velocity~$(2N)$. Wheeled kinematics: $(x,y,\theta,v,\delta)$~$(5N)$.}
\label{tab:obs_modes}
\begin{tabular}{l|ll}
    \toprule
    \textbf{Mode} & \textbf{UAV ($\dim$)} & \textbf{Wheeled ($\dim$)} \\
    \midrule
    \texttt{full}        & position + velocity + $\phi(\mathbf{b})$ ($4N+1$) & full kin.\ + $\phi(\mathbf{b})$ ($5N+1$) \\
    \texttt{no\_pos}     & $\phi(\mathbf{b})$ only ($1$) & $(\theta,v,\delta)$ per robot + $\phi(\mathbf{b})$ ($3N+1$) \\
    \texttt{no\_vis\_hist} & position + velocity ($4N$) & full kin.\ only ($5N$) \\
    \texttt{pos\_only}   & position only ($2N$) & position only ($2N$) \\
    \bottomrule
\end{tabular}
\end{table}

% ─────────────────────────────────────────────────────────────────────────────
\subsection{Setup}
\label{sec:exp:setup}
% ─────────────────────────────────────────────────────────────────────────────

We compare four on-policy (A2C \citep{mnih2016a2c}, ARS \citep{mania2018ars}, PPO \citep{schulman2017ppo}, TRPO \citep{schulman2015trust}) and two off-policy (CrossQ \citep{bhatt2024crossq}, TQC \citep{kuznetsov2020tqc}) algorithms via Stable-Baselines3 \citep{stable-baselines3}. Main experiments use: (A)~\textit{default hyperparameters} from Stable-Baselines3 and (B)~\textit{tuned hyperparameters} found by Optuna \citep{optuna_2019} (details in supplementary). For the \textbf{UAV} platform, we vary the number of robots $N \in \{2, 3, 4, 5\}$ and use 10 environment configurations (varied field shapes, starting positions, and target locations); for the \textbf{Wheeled} platform, we similarly vary $N \in \{2, 3, 4, 5\}$ and use 10 arena layouts with distinct obstacle arrangements. Every algorithm--configuration pair is trained for $2 \times 10^6$ timesteps with 3 seeds $(42, 123, 9999)$. The interquartile mean (IQM) \citep{agarwal2021deep} is reported as the primary metric. All experiments were conducted on an HPC cluster with GPUs: NVIDIA\textsuperscript{\textregistered} GeForce RTX 1080Ti, 2080Ti, 3090, L4, L40s, A4000, A4500, A6000; CPUs: Intel\textsuperscript{\textregistered} Xeon\textsuperscript{\textregistered} E5-2680v3, E5-2683v4, Gold 6130 (16--48 cores); and 4--16 GB RAM per node. Simulations were performed on a local desktop with AMD Ryzen\texttrademark\ Threadripper 1950X (16-core), 64 GB RAM, and NVIDIA\textsuperscript{\textregistered} GeForce RTX 2080. Each run took approximately 1--6 hours.

% ─── Table 1: UAV Main Results ───────────────────────────────────────────────
\begin{table*}[tbp]
\centering
\footnotesize
\caption{\small IQM of episodic reward for the \textbf{UAV} platform under default and tuned hyperparameters, over 3 seeds $\times$ 10 configurations $\times$ 50 episodes. $\dagger$~=~best per column. ($-$~=~placeholder; to be filled after experiments complete.)}
\label{tab:uav_main}
% \resizebox{\textwidth}{!}{
\begin{tabular}{l|cccc|cccc}
    \toprule
    \multirow{2}{*}{\textbf{Algorithm}}
        & \multicolumn{4}{c|}{\textbf{Default Hyperparameters}}
        & \multicolumn{4}{c}{\textbf{Tuned Hyperparameters}} \\
    \cmidrule(lr){2-9}
           & $N=2$ & $N=3$ & $N=4$ & $N=5$
           & $N=2$ & $N=3$ & $N=4$ & $N=5$ \\
    \midrule
    A2C    & $-$ & $-$ & $-$ & $-$ & $-$ & $-$ & $-$ & $-$ \\
    ARS    & $-$ & $-$ & $-$ & $-$ & $-$ & $-$ & $-$ & $-$ \\
    PPO    & $-$ & $-$ & $-$ & $-$ & $-$ & $-$ & $-$ & $-$ \\
    TQC    & $-$ & $-$ & $-$ & $-$ & $-$ & $-$ & $-$ & $-$ \\
    TRPO   & $-$ & $-$ & $-$ & $-$ & $-$ & $-$ & $-$ & $-$ \\
    CrossQ & $-$ & $-$ & $-$ & $-$ & $-$ & $-$ & $-$ & $-$ \\
    \bottomrule
\end{tabular}
\vspace{-6pt}
\end{table*}

% ─── Table 2: Wheeled Main Results ───────────────────────────────────────────
\begin{table*}[tbp]
\centering
\footnotesize
\caption{\small IQM of episodic reward for the \textbf{Wheeled} platform under default and tuned hyperparameters, over 3 seeds $\times$ 10 configurations $\times$ 50 episodes.
$\dagger$~=~best per column. ($-$~=~placeholder; to be filled after experiments complete.)}
\label{tab:whl_main}
% \resizebox{\textwidth}{!}{
\begin{tabular}{l|cccc|cccc}
    \toprule
    \multirow{2}{*}{\textbf{Algorithm}}
        & \multicolumn{4}{c|}{\textbf{Default Hyperparameters}}
        & \multicolumn{4}{c}{\textbf{Tuned Hyperparameters}} \\
    \cmidrule(lr){2-9}
           & $N=2$ & $N=3$ & $N=4$ & $N=5$
           & $N=2$ & $N=3$ & $N=4$ & $N=5$ \\
    \midrule
    A2C    & $-$ & $-$ & $-$ & $-$ & $-$ & $-$ & $-$ & $-$ \\
    ARS    & $-$ & $-$ & $-$ & $-$ & $-$ & $-$ & $-$ & $-$ \\
    PPO    & $-$ & $-$ & $-$ & $-$ & $-$ & $-$ & $-$ & $-$ \\
    TQC    & $-$ & $-$ & $-$ & $-$ & $-$ & $-$ & $-$ & $-$ \\
    TRPO   & $-$ & $-$ & $-$ & $-$ & $-$ & $-$ & $-$ & $-$ \\
    CrossQ & $-$ & $-$ & $-$ & $-$ & $-$ & $-$ & $-$ & $-$ \\
    \bottomrule
\end{tabular}
\vspace{-6pt}
\end{table*}

% ─────────────────────────────────────────────────────────────────────────────
\subsection{Main Results}
\label{sec:exp:main}
% ─────────────────────────────────────────────────────────────────────────────

Tables~\ref{tab:uav_main} and~\ref{tab:whl_main} compare all six algorithms under default and Optuna-tuned hyperparameters, evaluated with 50 episodes per trained model. \textit{(Results and analysis to be added once experiments are complete.)} For reference, an episode is considered successful when all $m$ targets are visited before the step limit, which activates the $+r_L$ terminal bonus in Eq.~\eqref{eq:r_term}.

% ─── Figures (learning curves) ───────────────────────────────────────────────
% \begin{figure}[tbp]
%     \centering
%     \includegraphics[width=\linewidth]{figures/combined_learning_curves.png}
%     \caption{Learning curves --- Default (top row) vs.\ tuned (bottom row) hyperparameters. Left panels: UAV platform. Right panels: Wheeled platform. Dotted lines indicate the reward threshold for a successful policy.}
%     \label{fig:curves}
%     \vspace{-12pt}
% \end{figure}

% ─────────────────────────────────────────────────────────────────────────────
\subsection{Ablation Study}
\label{sec:exp:ablation}
% ─────────────────────────────────────────────────────────────────────────────

We ablate reward terms, observation spaces, and uncertainty parameters using CrossQ ($N=3$, 3 seeds, 10 configurations, 50 episodes) on both platforms. Table~\ref{tab:ablations_and_dr} summarizes the results.

\textbf{Reward ablations} compare the IQM of episode lengths, where episodes terminate either on complete target coverage (all $m$ targets visited) or on a collision. Removing $\rterm$ prevents the policy from learning collision avoidance and coverage incentives. Removing $\rpath$ eliminates energy and path-efficiency penalties; because the discount factor $\gamma$ still implicitly penalizes slow episodes, the impact is expected to be moderate.

\textbf{Observation ablations} evaluate the effect of withholding kinematic information (\texttt{no\_pos}), visit history (\texttt{no\_vis\_hist}), or both (retaining \texttt{pos\_only}). We hypothesize that removing position information is catastrophic since robots cannot navigate, and that withholding visit history has a milder effect as robots can re-discover targets.

\textbf{Uncertainty ablations} isolate the contributions of wind noise (\texttt{wind\_only}), actuation noise (\texttt{act\_only}), neither (\texttt{deterministic}), or both (\texttt{full}). This identifies which noise source is the primary driver of performance degradation.

% ─── Table 3: Ablations + DR ─────────────────────────────────────────────────
\begin{table*}[tbp]
\centering
\small
\caption{\small Ablation and domain randomization results for $N=3$ using CrossQ, over 3 seeds $\times$ 10 configurations $\times$ 50 episodes. Left three blocks: ablation experiments. Right block: DR experiments. $\dagger$~=~best per block/column. UAV~(U) and Wheeled~(W) results shown per row.}
\label{tab:ablations_and_dr}
\resizebox{\textwidth}{!}{%
\begin{tabular}{lcc|lcc|lcc|lccc}
    \toprule
    \multicolumn{3}{c|}{\textbf{Reward ablation}}
    & \multicolumn{3}{c|}{\textbf{Observation ablation}}
    & \multicolumn{3}{c|}{\textbf{Uncertainty}}
    & \multicolumn{4}{c}{\textbf{Domain Randomization}} \\
    \midrule
    Condition & \makecell{Ep.\ len \\ IQM (U)} & \makecell{Ep.\ len \\ IQM (W)}
    & Condition & \makecell{Reward \\ IQM (U)} & \makecell{Reward \\ IQM (W)}
    & Condition & \makecell{Reward \\ IQM (U)} & \makecell{Reward \\ IQM (W)}
    & Condition & \makecell{In-dist.\ \\ IQM (U)} & \makecell{OOD \\ IQM (U)} & CVaR$_{0.1}$ \\
    \midrule
    Full          & $-$ & $-$ & Full          & $-$ & $-$ & Full          & $-$ & $-$ & No DR   & $-$ & $-$ & $-$ \\
    No $\rterm$   & $-$ & $-$ & No pos.       & $-$ & $-$ & Wind only     & $-$ & $-$ & Wind DR & $-$ & $-$ & $-$ \\
    No $\rpath$   & $-$ & $-$ & No visit hist.& $-$ & $-$ & Act.\ only    & $-$ & $-$ & Full DR & $-$ & $-$ & $-$ \\
                  &     &     & Pos.\ only    & $-$ & $-$ & Deterministic & $-$ & $-$ &         &     &     &     \\
    \bottomrule
\end{tabular}}
\vspace{-12pt}
\end{table*}

% ─────────────────────────────────────────────────────────────────────────────
\subsection{Domain Randomization}
\label{sec:exp:dr}
% ─────────────────────────────────────────────────────────────────────────────

We implement domain randomization (DR) \citep{tobin2017domain, peng2018sim2real} by re-sampling physical parameters per training episode under three conditions: (A)~No DR, (B)~Wind DR (randomize wind speed and direction uniformly over $[0, 1] \text{m/s}$), and (C)~Full DR (additionally randomize actuation noise std, robot mass/thrust for UAV, or max speed/steer for wheeled). We evaluate on in-distribution ($\text{wind} \in [0, 0.5] \text{m/s}$) and OOD ($\text{wind} \in (0.5, 2.0] \text{m/s}$) test conditions (right block of Table~\ref{tab:ablations_and_dr}). We additionally evaluate the conditional value-at-risk CVaR$_{0.1}$ \citep{tamar2015cvar} to measure worst-case robustness. We hypothesize that wind DR substantially improves OOD performance while full DR further reduces tail risk at some cost to mean performance.
\vspace{-4pt}

% ─────────────────────────────────────────────────────────────────────────────
\subsection{Simulation and Sim-to-Real Analysis}
\label{sec:exp:s2r}
% ─────────────────────────────────────────────────────────────────────────────

\begin{wraptable}{r}{7cm}
\vspace{-1.8em}
\centering
\caption{\small Sim-to-real observation-gap analysis (CrossQ, UAV, $N = 3$, configuration~1, 50 eval eps). $\Delta$IQM: relative change vs.\ no-perturbation baseline.}
\label{tab:obs_gap}
\setlength{\tabcolsep}{4pt}
\small
\begin{tabular}{lcc}
    \toprule
    \textbf{Perturbation} & \textbf{IQM} & $\Delta$\textbf{IQM (\%)} \\
    \midrule
    None (baseline)                      & $-$ & $0.0$   \\
    GPS noise ($\sigma = 0.05$ m)     & $-$ & $-$     \\
    Wind latency (3 steps)               & $-$ & $-$     \\
    Comm.\ loss (10\%)                   & $-$ & $-$     \\
    All combined                         & $-$ & $-$     \\
    \bottomrule
\end{tabular}
\vspace{-0.8em}
\end{wraptable}

We deploy trained UAV policies in a 3-D CoppeliaSim~\citep{coppeliaSim2013} agricultural scene with PyBullet~\citep{coumans2021pybullet} physics connected via ZeroMQ~\citep{ZeroMQ}. Table~\ref{tab:obs_gap} quantifies IQM degradation under inference-time sensor imperfections applied to the CrossQ policy trained with default hyperparameters ($N = 3$, configuration~1, full DR). Four perturbation conditions are evaluated: GPS position noise ($\sigma = 0.05$ m), wind estimate latency (3 steps), communication packet loss (10\%), and all combined. Wind estimate latency is expected to be the dominant transfer risk, as the UAV policy relies on timely wind information to compensate trajectory deviations. GPS position noise should have a smaller individual effect since the coverage reward is spatially smooth. Applying all perturbations simultaneously is expected to cause the largest collapse in performance, consistent with transfer gaps reported for UAV sim-to-real tasks \citep{peng2018sim2real, akkaya2019solving}. Fig.~\ref{fig:sim} illustrates navigation paths for two UAV environment configurations in the CoppeliaSim scene.

% \begin{figure}[tbp]
%     \centering
%     \subfloat[Config.~1 initial]{\includegraphics[width=0.247\textwidth, height=0.17\textwidth]{figures/simulations/env1_init.png}\label{fig:env1_init}}
%     \hfil
%     \subfloat[Config.~1 final]{\includegraphics[width=0.247\textwidth, height=0.17\textwidth]{figures/simulations/env1_v1_best.png}\label{fig:env1_best}}
%     \hfil
%     \subfloat[Config.~3 initial]{\includegraphics[width=0.247\textwidth, height=0.17\textwidth]{figures/simulations/env3_v1_init.png}\label{fig:env3_init}}
%     \hfil
%     \subfloat[Config.~3 final]{\includegraphics[width=0.247\textwidth, height=0.17\textwidth]{figures/simulations/env3_v1_best.png}\label{fig:env3_best}}
%     \caption{\small CoppeliaSim simulation of UAV navigation for configurations~1 (a, b) and~3 (c, d) with $N=3$ robots. Cyan markers indicate unvisited target locations; dark grey markers indicate visited targets; colored lines trace robot navigation paths.}
%     \label{fig:sim}
%     \vspace{-12pt}
% \end{figure}

\section{Conclusion}
\label{sec:con}

We introduced a dynamics-agnostic framework for multi-robot navigation and a dual-platform stochastic benchmark built on it. By formulating the cooperative target-coverage task using a general robot model and a unified free-workspace representation, the framework accommodates arbitrary motion dynamics and workspaces with or without static obstacles. We proved the equivalence between the constrained navigation optimization problem and the POMDP we construct, providing a formal guarantee linking the POMDP reward to the underlying constrained navigation objective. Instantiating the framework with UAV (Newtonian point-mass) and wheeled-robot (bicycle, non-holonomic) platforms under a shared regime of wind disturbances, actuation noise, and observation noise demonstrates its generality across robot types. Our comprehensive experimental protocol---covering six algorithms, two platforms, three seeds, and four categories of ablation---offers a standardized testbed for evaluating robustness to stochastic disturbances and variable motion dynamics.

Several directions remain open for future work. First, extending the framework to heterogeneous fleets (mixed UAV and wheeled teams) would test coordination across platforms with fundamentally different motion models. Second, incorporating robot failure or disconnection mid-episode would require policies to redistribute tasks dynamically among remaining active agents. Third, replacing the centralized observation model with a decentralized one---where each robot observes only neighbors within a communication range---would push toward more scalable and realistic deployments. Finally, physical deployment of the learned policies on real UAV and ground robot hardware would provide the ultimate validation of the sim-to-real transfer properties studied in this work.


%===============================================================================
\clearpage
% The acknowledgments are automatically included only in the final and preprint versions of the paper.
\acknowledgments{If a paper is accepted, the final camera-ready version will (and probably should) include acknowledgments. All acknowledgments go at the end of the paper, including thanks to reviewers who gave useful comments, to colleagues who contributed to the ideas, and to funding agencies and corporate sponsors that provided financial support.}
%===============================================================================

% no \bibliographystyle is required, since the corl style is automatically used.
\bibliography{ref}  % .bib

\end{document}
